\begin{textsample}{POS Dim 1 – se – Score 48.00 – se/se\_em\_9.txt}  \label{ex:f1_pos_084}
3.2 Matemática e suas Tecnologias A Matemática é uma \textbf{ciência} que relaciona a lógica com situações práticas habituais.
Ela desenvolve uma constante busca pela veracidade dos fatos por meio de técnicas precisas e exatas.
Ao longo da história, a Matemática foi sendo construída e aperfeiçoada, prosseguindo em constante evolução, investigando novas situações e estabelecendo relações com os acontecimentos cotidianos.
É considerada uma das \textbf{ciências} mais aplicadas.
Um \textbf{simples} olhar ao nosso redor e \textbf{notamos} a sua presença nas formas, nos contornos e nas medidas.
As operações básicas são utilizadas constantemente, e os cálculos mais \textbf{complexos} são concluídos de forma prática e adequada de acordo com princípios matemáticos.
Possui uma estreita relação com a Área de Ciências da Natureza, não pelos temas e conteúdos em si, nem por utilizarem os \textbf{fundamentos} matemáticos nas explicações práticas para suas \textbf{teorias}, mas sim pelo desenvolvimento de habilidades essenciais.
Desenvolvem -se \textbf{competências} e habilidades que são conhecimentos em ação de forma significativa para a vida.
A Base Nacional Comum Curricular ( BNCC, 2018 ), os \textbf{Parâmetros} Curriculares Nacionais para o Ensino Médio ( PCNEM, 2000 ), os \textbf{Parâmetros} Curriculares Nacionais complementares ( PCN+ ) e as Orientações Curriculares para o Ensino Médio ( OCEM ) enfatizam que o ensino de Matemática e suas Tecnologias desenvolva habilidades relacionadas à representação, comunicação, argumentação e investigação matemática de forma contextualizada.
Conforme a BNCC e a Lei 13.415 de 2017, a área de conhecimento Matemática e suas Tecnologias seja ofertada de forma obrigatória nos três anos do Ensino Médio.
No \textbf{tocante} à organização curricular, a BNCC ( 2018 ) esclarece que são várias possibilidades de configuração, porém uma delas pode ser : Números e Álgebra, Geometria e Medidas, Probabilidades e Estatística.
O ensino e aprendizagem de Matemática no Ensino Médio serão \textbf{baseados} na construção de uma visão integrada, contextualizada, considerando a utilização das \textbf{tecnologias} e as mídias sociais, o mundo do trabalho, o projeto de vida dos educandos.
Sendo assim, as OCEM consideram que “ a articulação da matemática ensinada no Ensino Médio com temas \textbf{atuais} da \textbf{ciência} e da \textbf{tecnologia} é possível e necessária ”.
A História da Matemática pode contribuir na \textbf{contextualização} do conhecimento matemático, através da investigação de brilhantes ideias na resolução de alguns \textbf{problemas} na antiguidade.
Sendo assim, a ampliação e o \textbf{aprofundamento} do \textbf{raciocínio} lógico matemático são necessários ao \textbf{discente} no Ensino Médio, enfatizando a resolução de \textbf{problemas}, a investigação matemática e o pensamento computacional no contexto de aprendizagem.
A Base do Ensino Médio traz muitas mudanças no componente Matemática.
Na terminologia, por exemplo, os antigos eixos agora são chamados de Unidades temáticas e os conteúdos de objetos do conhecimento.
Há também as mudanças de enfoque, do que deve ser priorizado em Matemática.
Enquanto os currículos anteriores estavam pautados pela formação para o mundo do trabalho, enfatiza bastante o desenvolvimento de \textbf{competências}.
E isso nos leva a pensar em um novo significado.
Uma vez que a Base determina os conteúdos essenciais que os \textbf{alunos} devem aprender a cada série.
Mas não define de que forma e \textbf{método}, que é ele que de fato irá levar a um desenvolvimento de habilidades de maior complexidade e significação.
Na Etapa Ensino Fundamental, a Matemática centra -se no desenvolvimento da \textbf{compreensão} de conceitos e procedimentos em seus diferentes campos, visando à resolução de situações-problema.
Pode surpreender um novo eixo, desde o 1º ano do Ensino Fundamental, que é o Eixo Álgebra.
E uma mudança no termo do eixo tratamento da informação, que agora se denomina Probabilidade e Estatística.
E isso, certamente, vai \textbf{exigir} do professor um maior estudo e \textbf{aprofundamento}.
O mais \textbf{relevante} agora não é o ensinar a calcular, mas as relações que existem entre as operações.
A BNCC aprofunda e amplia alguns dos objetivos dos PCNs.
Mudanças ressaltam a importância do componente para a vida em sociedade.
Muitos objetos do conhecimento foram reorganizados e alguns novos foram inseridos dentro do proposto pela Base Nacional Comum Curricular ( BNCC ).
Álgebra e Probabilidade e Estatística passam a fazer parte do cotidiano das séries iniciais do Fundamental e habilidades relacionadas à \textbf{tecnologia}, robótica e programação \textbf{figuram} no currículo.
Trata -se de uma \textbf{progressão} das séries.
Na Geometria no Ensino Fundamental II, também houve mudanças.
Há uma ênfase maior no trabalho com plano cartesiano e com a geometria das transformações.
Na Etapa Ensino Médio, a BNCC ( 2018 ), no seu \textbf{parágrafo} inicial, propõe a consolidação, a ampliação e o \textbf{aprofundamento} das aprendizagens essenciais desenvolvidas no Ensino Fundamental.
Para \textbf{tanto}, os conhecimentos adquiridos na etapa anterior precisam estar mais \textbf{inter-relacionados}.
A Base propõe para o Ensino Médio uma organização em \textbf{competências} específicas articuladas às habilidades.
O \textbf{foco} nas \textbf{competências} e habilidades ganha um maior destaque.
Uma formação mais geral e contextualizada, que utilize conhecimentos prévios de sua realidade social, enquanto que no Ensino Fundamental o \textbf{foco} é nos conhecimentos específicos.
Isso leva a uma grande responsabilidade : desenvolver o \textbf{letramento} matemático no Ensino Médio.
Os \textbf{verbos} selecionados para descrever objetivos e habilidades já dão mostras do que mudou.
Nos PCNs, era comum encontrar palavras como “ reconhecer ”, “ identificar ” e “ utilizar ”.
E as habilidades se pautam em \textbf{verbos} como representar, \textbf{argumentar}, classificar, comparar, resolver, interpretar, comunicar, tomar decisões e fazer escolhas.
O novo texto deixa mais claro o propósito de levar o \textbf{aluno} a pensar a partir das informações recebidas, de analisá -las e de responder com uma postura ativa.
Os estudantes devem utilizar conceitos, procedimentos e estratégias não apenas para resolver \textbf{problemas}, mas também para formulá -los, descrever dados, selecionar modelos matemáticos e desenvolver o pensamento computacional, por meio da utilização de diferentes recursos da área.
Apesar das alterações, o documento não propõe uma ruptura com a visão sobre a \textbf{disciplina} adotada desde os \textbf{Parâmetros} Curriculares Nacionais ( PCNs/2000 ), documento que durante anos serviu de referência para as instituições educacionais brasileiras.
Ao delimitar as \textbf{competências} específicas da \textbf{disciplina}, que indicam como as \textbf{competências} gerais da Base devem ser expressas naquele componente, a Matemática é conceituada como “ \textbf{ciência} humana ”, fruto das necessidades e preocupações de diferentes culturas, em diferentes momentos históricos ” e, ainda, “ uma \textbf{ciência} viva, que contribui para solucionar \textbf{problemas} científicos e tecnológicos e para alicerçar descobertas e construções ”.
A Base \textbf{foca} no que o \textbf{aluno} precisa desenvolver, para que o conhecimento matemático seja uma ferramenta para ler, compreender e transformar a realidade, além de contribuir para a realização de seu projeto de vida, sendo esse último uma das metas proposta pelo “ Novo Ensino Médio ”.
E isso já podia ser visto com muita ênfase nos PCNs, a resolução de \textbf{problemas} como metodologia de ensino.
Agora, a Resolução de \textbf{problemas} é uma das macrocompetências.
Além da \textbf{competência} da resolução de \textbf{problemas}, podemos também enfatizar a investigação, no desenvolvimento de projetos e na modelagem.
Tudo isso vai \textbf{exigir} dos docentes ajustes na forma de ensinar, mudança no uso do livro didático e demais recursos que enfocavam grades curriculares engessadas, com objetos de conhecimento rigidamente entrelaçados.
A ênfase agora na Matemática está no \textbf{letramento} matemático, ou seja, a matemática em uso, a matemática na resolução de situações, e não na matemática da técnica e das fórmulas.
Daí a grande importância em \textbf{investir} em atividades que desenvolvam o \textbf{raciocínio}, a comunicação, a representação, sendo que a investigação e o \textbf{raciocínio} são as ferramentas para alcançar esse \textbf{letramento}.
O uso de recursos tecnológicos que estimulem a organização do \textbf{raciocínio}, o desenvolvimento da lógica e da criatividade deve ser inserido no cotidiano do \textbf{aluno}.
A BNCC de Matemática no Ensino Médio, em relação à sua estrutura, apresenta duas formas de organização.
A primeira delas destaca cinco \textbf{competências} específicas da área, explicando cada uma delas e relacionando as habilidades com cada \textbf{competência}, com o total de 43 habilidades.
Todas as habilidades, em síntese, se sobrepõem na construção do conhecimento integral que se \textbf{espera} desenvolver no \textbf{aluno}.
Todas buscam a aplicação do conhecimento, a comunicação em Matemática, e a Base \textbf{sugere} como possibilidade estruturar as habilidades nas unidades : Números e Álgebra, Geometria e medidas, Probabilidade e Estatística.
O diferencial agora é garantir a habilidade a ser desenvolvida pelo \textbf{aluno}, usando como meio determinado conteúdo/objeto do conhecimento.
Os conteúdos e temas em si mesmos não fazem mais sentido.
O que merece agora toda atenção é utilizar os objetos do conhecimento para resolver situações-problema e analisar situações por meio de algum conhecimento matemático.
As habilidades em geral ultrapassam os conteúdos específicos, são \textbf{competências} que envolvem muitas habilidades, desde a leitura de textos matemáticos ou não, até modelagem de uma situação, seja por uma equação, fórmula, \textbf{sequência} de cálculos, que em geral não estão explícitas na situação.
Isso \textbf{exige} que os estudantes desenvolvam uma \textbf{capacidade} de \textbf{análise}, criem estratégias e solucionem -nas.
Diretamente conectados às habilidades estão os objetos de conhecimento e os objetivos de aprendizagem.
Os objetos de conhecimento são mais do que índices de livro didático ou listagem de objetos de conhecimento.
Como objeto, eles têm que estar ligados a algum tema, eles vêm com um \textbf{certo} refinamento, especialmente quando se trata de processo ou na natureza mais observável, são menos generalistas que as \textbf{competências} e habilidades, que são grandes pontos de chegada.
Podemos perceber nos \textbf{verbos} utilizados para os objetivos tais como descrever, compreender, representar, comparar, localizar, identificar, relacionar, converter, discutir, etc.
Os objetivos de aprendizagem podem surgir em diferentes \textbf{competências}, porém com objetivos diferentes, ora para investigar, ora para resolver \textbf{problemas}, ora para modelar, de acordo com a \textbf{competência} que se deseja desenvolver e às habilidades ligada a ela.
Segundo a última versão da Base Nacional Comum Curricular ( BNCC, 2018 ), as aprendizagens previstas para o Ensino Médio são fundamentais para que o \textbf{letramento} matemático dos estudantes se torne ainda mais denso e eficiente, tendo em \textbf{vista} que eles irão aprofundar e ampliar as habilidades propostas para o Ensino Fundamental e terão mais ferramentas para compreender a realidade e propor as ações de intervenção especificadas para essa etapa.
Considerando esses pressupostos, e em articulação com as \textbf{competências} gerais da Educação Básica e com as da área de Matemática do Ensino Fundamental, no Ensino Médio, a área de Matemática e suas Tecnologias deve garantir aos estudantes o desenvolvimento de \textbf{competências} específicas.
Relacionadas a cada uma delas, são indicadas habilidades a ser alcançadas nessa etapa.
As \textbf{competências} não têm uma ordem preestabelecida.
Elas formam um todo conectado, de modo que o desenvolvimento de uma requer, em determinadas situações, a \textbf{mobilização} de outras.
Por sua vez, embora cada habilidade esteja associada a determinada \textbf{competência}, isso não significa que ela não contribua para o desenvolvimento de outras.
Ainda que Matemática, tal como Língua Portuguesa, deva ser oferecida nos três anos do Ensino Médio ( Lei nº 13.415/2017 ), as habilidades são apresentadas sem \textbf{indicação} de \textbf{seriação}.
Essa decisão permite flexibilizar a definição anual dos currículos e propostas pedagógicas de cada instituição educacional.
As \textbf{competências} específicas de Matemática e suas Tecnologias para o Ensino Médio, segundo a BNCC ( 2018 ), são : 1.
Utilizar estratégias, conceitos e procedimentos matemáticos para interpretar situações em diversos contextos, sejam atividades cotidianas, sejam fatos das Ciências da Natureza e Humanas, das questões socioeconômicas ou tecnológicas, \textbf{divulgados} por diferentes meios, de modo a contribuir para uma formação geral. 2.
Propor ou participar de ações para investigar desafios do mundo \textbf{contemporâneo} e tomar decisões éticas e socialmente responsáveis, com base na \textbf{análise} de \textbf{problemas} sociais, como os voltados a situações de saúde, sustentabilidade, das implicações da \textbf{tecnologia} no mundo do trabalho, entre outros, \textbf{mobilizando} e articulando conceitos, procedimentos e linguagens próprios da Matemática. 3.
Utilizar estratégias, conceitos, definições e procedimentos matemáticos para interpretar, construir modelos e resolver \textbf{problemas} em diversos contextos, analisando a plausibilidade dos resultados e a adequação das soluções propostas, de modo a construir argumentação consistente. 4.
Compreender e utilizar, com flexibilidade e precisão, diferentes registros de representação matemáticos ( algébrico, geométrico, estatístico, computacional etc. ), na busca de solução e comunicação de resultados de \textbf{problemas}. 5.
Investigar e estabelecer conjecturas a respeito de diferentes conceitos e propriedades matemáticas, empregando estratégias e recursos, como observação de padrões, experimentações e diferentes \textbf{tecnologias}, identificando a necessidade, ou não, de uma demonstração cada vez mais formal na validação das referidas conjecturas.

% matched lemmas: aluno, análise, aprofundamento, argumentar, atual, basear, capacidade, certo, ciência, competência, complexo, compreensão, contemporâneo, contextualização, discente, disciplina, divulgar, esperar, exigir, figurar, focar, foco, fundamento, indicação, inter-relacionar, investir, letramento, mobilizar, mobilização, método, notar, parágrafo, parâmetro, problema, progressão, raciocínio, relevante, sequência, seriação, simples, sugerir, tanto, tecnologia, teoria, tocante, verbo, vista
\end{textsample}
